%%%%%%%%%%%%%%%%%%%%%%%%%%%%%%%%%%%%%%%%%%%%%%%%%%%%%%%%%%%%%%%%%%%%%%%%
\section{Area versus number of squares}\label{sec:area}
%%%%%%%%%%%%%%%%%%%%%%%%%%%%%%%%%%%%%%%%%%%%%%%%%%%%%%%%%%%%%%%%%%%%%%%%

The field also defines an area. Let $\nu$ be the Gaussian multiplicative chaos of $h$ with
parameter $\gamma$, the measure $e^{\gamma h(z)-\gamma^2\E[h(z)^2]/2}dz$ of
Section~\ref{sec:upper}; this is the Liouville quantum gravity area of the surface. A
square $S$ has area roughly
\[
\nu(S)\approx e^{\gamma h_S-\gamma^2\log(1/|S|)/2}|S|^2=e^{\gamma h_S}|S|^{2+\gamma^2/2}
=e^{\gamma h_S}|S|^{\gamma Q},
\]
using $\gamma Q=2+\gamma^2/2$. Raising the rule~\eqref{eq:rule} to the power $\gamma Q$ shows that a
square is cut exactly when this rough area exceeds $\varepsilon^{\gamma Q}$. So one might guess that
$N$ is about $\nu(S_0)/\varepsilon^{\gamma Q}$. By a theorem of Wong~\cite{Wong}, $\Prob(\nu(S_0)>t)$
is asymptotic to a constant times $t^{-4/\gamma^2}$, the exponent of Theorem~\ref{thm:tail}. Yet
the guess is false in a precise sense.

\begin{proposition}\label{prop:area}
For every $Q>2$ and fixed $\varepsilon>0$,
\[
\liminf_{t\to\infty}\Prob\bigl(N=1\;\big|\;\nu(S_0)>t\bigr)>0 .
\]
\end{proposition}

\begin{proof}
Let $Y=h_{S_0}$ be the average over the unit square; it is a nondegenerate centered normal
variable, and $N=1$ exactly when $Y\le Q\log\varepsilon$, since then $S_0$ is not active. Write
$h(z)=Y\chi(z)+h^\perp(z)$ with $\chi(z)=\operatorname{Cov}(h(z),Y)/\operatorname{Var}(Y)$; the
residual field $h^\perp$ is independent of $Y$. On a closed subsquare $U$ in the interior of $S_0$,
$\chi$ is smooth and bounded, and $h^\perp$ still has covariance $-\log|z-w|$ plus a smooth
bounded function. So Wong's theorem applies to the chaos $\nu^\perp$ of $h^\perp$ and gives
that $t^{4/\gamma^2}\Prob(\nu^\perp(U)>t)$ has a positive limit.

The two chaos measures differ on $U$ by the factor $\exp(\gamma Y\chi(z)-\gamma^2\operatorname{Var}(Y)\chi(z)^2/2)$.
Choose a compact interval $\mathcal I$ below $Q\log\varepsilon$ with $\Prob(Y\in\mathcal I)>0$. For $Y\in\mathcal I$
this factor is at least a deterministic $d_0>0$ on $U$. By independence,
\[
\Prob\bigl(\nu(S_0)>t,\ N=1\bigr)\ge\Prob(Y\in\mathcal I)\,\Prob\bigl(\nu^\perp(U)>t/d_0\bigr),
\]
which is at least a positive constant times $t^{-4/\gamma^2}$, while $\Prob(\nu(S_0)>t)$ is at most a
constant times $t^{-4/\gamma^2}$, again by Wong's theorem. Dividing proves the claim.
\end{proof}

The rule~\eqref{eq:rule} sees the field only through averages over squares it has already made. If
the average over $S_0$ is low, the subdivision stops at once, however high the field is on fine
scales inside, while the area sees all scales. Given a large area, the proposition rules out that
$N$ is comparable to the area with conditional probability tending to one; it does not rule out a
comparison that holds with positive conditional probability. This is why the proof of
Theorem~\ref{thm:tail} works with the squares directly.
